%%%%% Logo após o chapter 2 inicial

Nesta pesquisa de Mestrado a dinâmica orbital e de atitude será simulada computacionalmente para resolver o problema de uma espaçonave tipo manipulador robótico em missão de \emph{rendezvous and docking/berthing} (RVD/B).
Esta espaçonave está inserida em uma missão de \emph{rendezvous and docking/berthing} (RVD/B) em órbita terrestre baixa (LEO, \emph{Low Earth Orbit}).
A formulação contempla:
(i) dinâmica orbital relativa, entre o \emph{chaser} e o \emph{target}, descrita pelas equações de Hill/Clohessy--Wiltshire para órbita circular;
(ii) a dinâmica de atitude do corpo completo e do manipulador robótico, modelados segundo a abordagem de corpos rígidos.

No subsistema orbital considera-se que a aproximação final é realizada ao longo do eixo radial (\emph{R-bar}) e que a transferência inicial entre órbitas circulares é obtida por meio da transferência de Hohmann, sem perturbações ambientais.  
Para o subsistema de atitude são analisadas as equações de movimento da espaçonave com base flutuante sujeita à ação dos torques de controle, bem como as equações de movimento das juntas do manipulador composto por quatro articulações revolutas e uma prismática (configuração 4R1P), responsáveis pelo alinhamento e extensão da ferramenta de captura.
